%=====================================================================
\chapter{Case Studies in Enterprise and Cloud Virtualization}\label{ch:casestudies}
%=====================================================================
\locator{6}{Part VI --- Emerging Hardware and Practice}

\begin{outcomes}
By the end of this chapter you will be able to:
\begin{tight}
  \item \textbf{Analyse} a virtualization decision using the arithmetic of the
        preceding chapters rather than the vendor's framing.
  \item \textbf{Identify} the measurement that would have changed each
        decision, and \textbf{say} when it should have been taken.
  \item \textbf{Explain} why a technically correct migration can still be a
        commercial failure.
  \item \textbf{Apply} the consolidation, overcommit, migration and cost models
        to a case you have not seen before.
  \item \textbf{Write} a recommendation that states its assumptions and the
        evidence that would overturn it.
\end{tight}
\end{outcomes}

\begin{prereq}
Everything. Each case below uses the arithmetic of a specific earlier chapter,
named where it arises, and the point of the chapter is that the arithmetic is
the same whatever the vendor calls the product.
\end{prereq}

\begin{keyterms}
consolidation ratio $\bullet$ right-sizing $\bullet$ lift and shift $\bullet$
refactoring $\bullet$ egress $\bullet$ reserved capacity $\bullet$ licence
model $\bullet$ blast radius $\bullet$ edge site $\bullet$ intermittent
connectivity $\bullet$ post-mortem $\bullet$ counterfactual
\end{keyterms}

\section{How to Read These}

Four cases. The figures are representative of real deployments rather than
drawn from one named organisation, and they are constructed so that the
arithmetic can be checked --- that is what they are for. Each follows the same
shape: what was decided, what happened, what the numbers said, and the single
measurement that would have changed the outcome.

\begin{alertbox}[Every number in this chapter is invented, and says so]
The cases are composites. They are built to be arithmetically consistent and to
exercise the models of earlier chapters, not to report on any particular
institution. Where a real, dated fact is used --- a licence change, a published
vulnerability --- it is cited as such in the chapter that introduced it.

\smallskip
Treat the figures as you would a worked example: reproduce them, change the
inputs, and see where the recommendation flips. The skill being taught is the
flipping point, not the answer.
\end{alertbox}

\section{Case 1: The Consolidation That Nearly Failed}

\begin{conceptbox}[The decision]
A university's computing department had 64 physical servers, bought one at a
time over nine years. It consolidated onto 5 hosts and a shared array. The
business case cited the arithmetic of \chref{what}: a 12.8:1 ratio, 85\% less
power, one estate to patch instead of 64.

\smallskip
Eleven weeks after the migration the estate was described by its own users as
``slower than the old servers''.
\end{conceptbox}

\textbf{What the numbers said.} Sizing had been done on \emph{allocated}
capacity, as \chref{cpumem} warns.

\[ \sigma_{\text{cpu}} = \frac{64 \text{ guests} \times 4 \text{ vCPU}}
   {5 \times 48 \text{ cores}} = \frac{256}{240} = 1.07 \]

which looks comfortable. But N+1 was not applied, and one host was taken for
maintenance in week 11:

\[ \sigma_{\text{cpu}}\big|_{N+1} = \frac{256}{4 \times 48} = \frac{256}{192}
   = 1.33 \]

That alone is survivable. What was not survivable was memory. The 64 old
servers had 32 GB each; the guests were given 32 GB each because that is what
they had before, and nobody measured what they used:

\[ \text{allocated} = 64 \times 32 = 2{,}048 \text{ GB} \qquad
   \text{installed} = 5 \times 384 = 1{,}920 \text{ GB} \]

\[ \sigma_{\text{mem}} = \frac{2{,}048}{1{,}920 \times 0.85} = 1.26 \]

With all five hosts this was absorbed by page sharing and ballooning
(\chref{cpumem}). With four hosts it was not:

\[ \sigma_{\text{mem}}\big|_{N+1} = \frac{2{,}048}{4 \times 384 \times 0.85}
   = \frac{2{,}048}{1{,}306} = 1.57 \]

The hosts reached the bottom of the reclamation ladder and began swapping.
\%RDY was fine throughout, which is why the first week of investigation looked
at the wrong resource.

\textbf{The measurement that would have changed it.} Active memory, before
migration. Measured afterwards it was 9.4 GB per workload:

\[ 64 \times 9.4 = 602 \text{ GB} \]

against 1{,}920 GB installed --- the estate needed a third of what was bought,
and the guests should have been given 12 GB, not 32.

\textbf{What was done.} Right-sizing, over three weeks, guest by guest.
$\sigma_{\text{mem}}$ fell to 0.59 and $\sigma_{\text{mem}}|_{N+1}$ to 0.74. No
hardware was added.

\begin{conceptbox}[The lesson, stated generally]
\emph{Migrating a server is not sizing a guest.} The old machine's
specification was a purchasing decision from up to nine years earlier, made
against a different workload, with headroom for growth that never arrived.
Carrying it across preserves every bad guess ever made and adds the
hypervisor's overhead on top.

\smallskip
Measure for a month before migrating; size to the 95th percentile with a stated
margin; and apply N+1 to \emph{every} ratio you quote, because the failure
happens on the day one host is away.
\end{conceptbox}

\section{Case 2: The Cloud Migration That Cost More}

\begin{conceptbox}[The decision]
A software house of 40 engineers moved its entire estate --- 120 guests,
development and production --- to a public cloud over four months. The
projection showed a 15\% saving. The first full month's invoice was 2.3 times
the projection, and the second was higher.
\end{conceptbox}

\textbf{Where the money went.} Four causes, in descending order, each traceable
to a chapter.

\rowcolors{2}{white}{lightbg}
\begin{longtable}{@{}L{3.7cm}L{2.2cm}L{4.3cm}L{4.5cm}@{}}
\toprule
\rowcolor{primary!15}
\textbf{Cause} & \textbf{Share of overrun} & \textbf{What happened} &
\textbf{Chapter} \\
\midrule
\endhead
Lift and shift without right-sizing & 44\% & 120 guests moved at their existing
vCPU and memory, which had never been measured & \chref{cpumem},
\chref{cloud} \\
On-demand pricing throughout & 31\% & No reservations, because nobody was
confident enough to commit. \chref{cloud}'s worked example: on demand was 61\%
more than owning & \chref{cloud} \\
Egress & 17\% & Nightly backups written \emph{out} of the cloud to the old
on-premises array, at \$0.09 per GB & \chref{cloud} \\
Development guests running continuously & 8\% & 44 non-production guests
billed $24 \times 7$ for work done $8 \times 5$ & \chref{cloud} \\
\bottomrule
\end{longtable}

\textbf{The arithmetic of the fix.} All four were addressed in six weeks, with
no architectural change.

\begin{tight}
  \item \textbf{Right-sizing.} Measured p95 across 120 guests; mean allocation
        fell from 6.2 vCPU and 21 GB to 2.1 vCPU and 7 GB. Compute cost
        $\times\,0.34$.
  \item \textbf{Reservations} on the 61 guests that were genuinely steady.
        \chref{cloud}'s figure for a three-year commitment is 55\%; a one-year
        commitment was taken instead at 34\%, because the business would not
        commit for three.
  \item \textbf{Backups kept in the cloud}, with only a weekly copy leaving.
        Egress fell by about 85\%.
  \item \textbf{Development guests stopped out of hours}, a scheduled task:
        $24 \times 7 = 168$ hours becomes $10 \times 5 = 50$, a 70\%
        reduction on those 44 guests.
\end{tight}

The invoice settled at about 0.78 of the on-premises baseline --- close to the
original projection, reached eight months late.

\begin{conceptbox}[The lesson: lift and shift is a migration strategy, not a cost strategy]
Lifting and shifting is a perfectly reasonable way to \emph{move}: it is fast,
it is low-risk, and it decouples the move from the redesign.

\smallskip
What it is not is a way to save money, and the reason is structural. On
premises, waste is invisible --- a guest with four times the memory it needs
costs nothing extra once the host is bought. In the cloud, every unit of waste
is a line on a monthly invoice. The cloud did not make this estate expensive;
it made it \emph{legible}, and what became legible was eleven years of
unmeasured allocation.

\smallskip
If a migration business case claims savings, ask which of the four items above
it has budgeted for. If the answer is none, the projection is wrong by roughly
the factor this case saw.
\end{conceptbox}

\section{Case 3: Edge Virtualization at Twenty-Three Sites}

\begin{conceptbox}[The decision]
A logistics company ran a small application at each of 23 depots: stock,
weighing, dispatch. Each site had two aging PCs under a desk. Connectivity to
head office was a consumer broadband line, and it failed somewhere in the
estate most weeks.

\smallskip
The proposal was to centralise everything into the data centre and give each
depot a thin client (\chref{desktop}).
\end{conceptbox}

\textbf{Why centralisation was rejected.} The arithmetic was not about cost but
about \chref{migration}'s availability vocabulary.

Depot operations must continue when the line is down. Measured over a year,
each site's link had about 99.2\% availability:

\[ (1 - 0.992) \times 8{,}760 = \mathbf{70 \text{ hours a year}} \]

Seventy hours a year per site of \emph{no dispatch at all}, against the current
arrangement where a line failure meant only that head office reporting was
delayed. Centralisation would have improved manageability and made the business
worse.

\textbf{What was built instead.} One small two-node cluster per depot, running
three guests: the application, a local database and a management agent.
Specifically:

\begin{tight}
  \item \textbf{Immutable images} (\chref{management}). Depots have no
        technical staff, so nothing is configured on site. A new version is a
        new image; rollback is the previous image.
  \item \textbf{Containers, not VMs, for the application} (\chref{containers}),
        because the hardware budget per site was small and the start-up speed
        made unattended recovery practical.
  \item \textbf{Local persistence with asynchronous replication} to head
        office. RPO of 15 minutes, accepted explicitly: a depot losing 15
        minutes of dispatch records in a total site failure was judged
        tolerable, and synchronous replication over a failing consumer line was
        not possible anyway.
  \item \textbf{No live migration} (\chref{migration}). Two nodes, so quorum is
        impossible; a witness ran at head office, and when the line was down
        the site operated in a documented degraded mode with HA disabled
        rather than risking the split brain of \chref{migration}.
\end{tight}

\begin{conceptbox}[The lesson: the constraint is rarely the one in the proposal]
The proposal was framed as a management problem --- 46 unmanaged PCs --- and
the answer was framed as centralisation. The actual constraint was a 99.2\%
network link, which nobody had measured until someone asked for the number.

\smallskip
Edge deployments invert several assumptions this book has relied on: bandwidth
is scarce and unreliable, there is no one on site, physical security is weak,
and the hardware is whatever fits under a desk. The techniques still apply ---
images, containers, replication, quorum --- but the \emph{weightings} change
completely.

\smallskip
The transferable question: \emph{what does this site have to keep doing when
the link is down?} If the answer is ``nothing'', centralise. If it is
``everything'', the compute follows the work.
\end{conceptbox}

\section{Case 4: A Licence Change, and What It Revealed}

\begin{conceptbox}[The decision]
A hospital trust ran 340 guests on 14 hosts under vSphere. At renewal, a change
in the vendor's packaging --- bundled subscription rather than perpetual
per-socket licensing, as the industry saw after 2024 (\chref{platforms}) ---
raised the quoted renewal by a factor of roughly four.

\smallskip
Three options were considered: pay, migrate to another platform, or move to a
public cloud.
\end{conceptbox}

\textbf{What the comparison looked like.} Over three years, approximately:

\rowcolors{2}{white}{lightbg}
\begin{longtable}{@{}L{4.0cm}L{2.6cm}L{3.2cm}L{4.6cm}@{}}
\toprule
\rowcolor{primary!15}
\textbf{Option} & \textbf{Licence} & \textbf{Migration effort} &
\textbf{Risk} \\
\midrule
\endhead
Stay on vSphere & High & None & Exposed to the next repricing \\
Migrate to Proxmox or oVirt & Near zero & 340 guests, 9--14 months, two
engineers & Staff retraining; tooling rebuilt; clinical systems vendors may not
support it \\
Public cloud & N/A & Largest; also a data-residency review & Egress; ongoing
cost; regulatory \\
\bottomrule
\end{longtable}

\textbf{The finding that decided it.} Not cost. When the trust catalogued what
would have to be rebuilt, it discovered that \textbf{nobody could enumerate what
the estate did}. There was no inventory mapping guests to clinical services, no
record of which of the 340 were still needed, and 61 guests whose owner could
not be identified at all.

The migration could not be planned because the thing to be migrated was not
known.

\textbf{What was done.} A twelve-month discovery and rationalisation programme
first: inventory, ownership, dependency mapping, decommissioning. It retired 78
guests --- 23\% of the estate --- that nothing used. Only then was the platform
decision taken, and by then the licence quantity had fallen with the guest
count.

\begin{conceptbox}[The lesson: lock-in is not technical]
Every chapter in this book would say the migration is feasible. The hypervisors
are comparable (\chref{platforms}), the disk formats convert
(\chref{storage}), the networking is reproducible (\chref{network}).

\smallskip
What made it hard was institutional knowledge that had never been written down.
That is the real lock-in, and it is not specific to any vendor: it accumulates
in any estate that grows without an inventory, and it is what makes every
platform change cost three times its estimate.

\smallskip
The cheapest insurance against a repricing you cannot predict is an inventory
you maintain. It costs little, it is useful every day for other reasons, and it
converts a crisis into a project.
\end{conceptbox}

\section{What the Four Have in Common}

\rowcolors{2}{white}{lightbg}
\begin{longtable}{@{}L{2.0cm}L{4.4cm}L{4.3cm}L{4.1cm}@{}}
\toprule
\rowcolor{primary!15}
\textbf{Case} & \textbf{Framed as} & \textbf{Actually about} &
\textbf{Missing measurement} \\
\midrule
\endhead
1 & Consolidation ratio & N+1 memory under maintenance & Active memory, before
migrating \\
2 & Cloud versus on premises & Eleven years of unmeasured allocation becoming
visible & p95 utilisation per guest \\
3 & Managing 46 unmanaged PCs & What must keep working when the link fails &
Link availability per site \\
4 & A licence renewal & Not knowing what the estate does & An inventory with
owners \\
\bottomrule
\end{longtable}

\begin{conceptbox}[The pattern]
In all four, the decision was framed as a technology choice and settled by a
measurement nobody had taken. In none of them was the chosen technology wrong.

\smallskip
That is worth carrying into \chref{project} and into professional practice. The
engineering in this book is the easy part: hypervisors work, containers work,
migration works. The difficulty is almost always in knowing what you have, what
it uses and what it must keep doing --- and those are questions of measurement
and record-keeping, not of architecture.

\smallskip
When you are handed a proposal, the most useful question is rarely ``is this
the right technology?''. It is \emph{``what number would change this
decision, and has anyone taken it?''}
\end{conceptbox}

\begin{pitfall}
\begin{tight}
  \item \textbf{Quoting a consolidation ratio without N+1.} Case 1 failed on
        the day one host went away for maintenance.
  \item \textbf{Carrying old specifications into new guests.} You preserve
        every purchasing guess ever made and add overhead.
  \item \textbf{Expecting lift and shift to save money.} It moves the estate;
        it does not fix it, and the cloud bills for what on premises concealed.
  \item \textbf{Centralising without measuring the link.} Availability at the
        edge is a property of the worst site, not the average one.
  \item \textbf{Treating a platform migration as a technical project.} The cost
        is in the knowledge that was never written down.
  \item \textbf{Reading a case study as a recommendation.} Change one input and
        the answer changes; the skill is knowing which input.
\end{tight}
\end{pitfall}

\begin{lab}[Analysing a Case You Have Not Seen]
No commands. Work in pairs, 45 minutes, and write a one-page recommendation ---
this is deliberately the shape of the examination's applied questions and of
\chref{project}'s report.

\smallskip
\textbf{The brief.} A college runs 90 guests on 4 hosts (32 cores, 256 GB
each), bought three years ago, under a platform whose support expires in seven
months. Measured over the last term:

\begin{tight}
  \item allocated: 410 vCPU, 1{,}740 GB
  \item active: 198 GB total, p95; \%RDY p95 11\%
  \item storage: 24 TB provisioned on an 11 TB thin pool, 7.9 TB written,
        growing 180 GB a month
  \item 14 guests have no identified owner
  \item the finance office requires three-year costings
\end{tight}

\textbf{Produce:}
\begin{enumerate}[leftmargin=2.2em, itemsep=3pt, topsep=4pt]
  \item $\sigma_{\text{cpu}}$ and $\sigma_{\text{mem}}$, with and without N+1.
  \item The storage over-subscription ratio and the time to pool exhaustion
        (\chref{storage}).
  \item The single most urgent finding, with the arithmetic that supports it.
  \item Three options with three-year costings, using \chref{cloud}'s method
        for the cloud one.
  \item A recommendation, the assumptions it rests on, and \textbf{the one
        measurement that would overturn it}.
\end{enumerate}

\textbf{Marking yourselves.} A good answer identifies that \%RDY of 11\% and
$\sigma_{\text{cpu}}$ of 4.0 mean the estate is already contended; that memory
is wildly over-allocated at 1{,}740 GB against 198 GB active; that the thin
pool has about 17 months at the current rate but no N+1 storage headroom; and
that the 14 unowned guests are case 4 in miniature. A weak answer compares
products.
\end{lab}

\begin{checkpoint}
\begin{tightnum}
  \item In case 1, $\sigma_{\text{mem}}$ was 1.26 with five hosts and 1.57 with
        four. Explain why the first was survivable and the second was not,
        naming the mechanism.
  \item Case 2's estate was not made more expensive by the cloud. State in one
        sentence what the cloud actually did, and name the two items that
        together accounted for 75\% of the overrun.
  \item In case 3, why did a two-node cluster at each depot need a witness
        elsewhere, and what was done when the link to it failed?
\end{tightnum}
\end{checkpoint}

\begin{chaptersummary}
\begin{tight}
  \item Case 1: a consolidation sized on allocated rather than active memory
        failed the day a host was taken for maintenance. Quote every ratio with
        N+1.
  \item Case 2: a cloud migration cost 2.3 times its projection through lift
        and shift, on-demand pricing, egress and always-on development guests.
        The cloud made existing waste legible rather than creating it.
  \item Case 3: centralising 23 depots would have improved management and cost
        70 hours a year of operations per site. The constraint was link
        availability, which nobody had measured.
  \item Case 4: a licence repricing was blocked not by technology but by the
        absence of an inventory. Lock-in is institutional, not technical.
  \item In all four the decision was framed as a technology choice and settled
        by a missing measurement.
  \item The question to ask of any proposal: what number would change this, and
        has anyone taken it?
\end{tight}
\end{chaptersummary}

\begin{reviewq}
  \item \textbf{(MCQ)} In case 1, the resource that failed was: (a)~CPU;
        (b)~memory under N+1; (c)~storage capacity; (d)~network bandwidth.
  \item \textbf{(MCQ)} The largest single contributor to case 2's overrun was:
        (a)~egress; (b)~on-demand pricing; (c)~lift and shift without
        right-sizing; (d)~development guests running continuously.
  \item \textbf{(Short)} Explain why the cloud ``made waste legible'' and what
        that implies for a migration business case.
  \item \textbf{(Short)} State the question that decides whether an edge site
        should keep local compute, and apply it to a supermarket till system.
  \item \textbf{(Short)} Explain what case 4 means by lock-in being
        institutional rather than technical.
  \item \textbf{(Applied)} Work the lab's brief and submit the one-page
        recommendation.
  \item \textbf{(Applied)} Choose any one of the four cases and write the
        counterfactual: what would have happened had the missing measurement
        been taken six months earlier? Be specific about what would have been
        decided differently and what it would have cost.
\end{reviewq}
